% Adapted from the original web post:
%   "CADI: a prescriptive path for slop-free AI in content work"
%   https://www.allaboutken.com/posts/20260609-iterative-pattern-framework/
% Restructured into paper form as a proof of concept for depositing on
% Zenodo (user: khawkins98) with an ORCID-linked author record.
\documentclass[11pt]{article}

\usepackage[margin=1in]{geometry}
\usepackage{hyperref}
\usepackage{enumitem}
\usepackage{parskip}

\title{CADI: A Prescriptive Framework for Evidence-Based AI-Assisted Content Work}
\author{%
  Ken Hawkins\thanks{Independent Researcher.} \\
  \href{https://orcid.org/0009-0009-6728-2237}{ORCID: 0009-0009-6728-2237}
}
\date{June 2026}

\begin{document}
\maketitle

\begin{abstract}
Most AI-writing assistance starts cold on every task: no memory of what
worked last time, no documented picture of an organisation's voice or
standards. The common responses --- better prompting, a static style
guide, a bolted-on agent --- each capture part of the answer but not
the order or the combination. This paper presents CADI (Collect,
Analyze, Document, Iterate), a four-phase, evidence-based loop for
deriving AI writing guidance from an organisation's own best content
and revising that guidance on measured outcomes rather than intuition.
The framework is illustrated with a case study from UNDRR (the UN
Office for Disaster Risk Reduction) and situated against related work
in voice-profile tooling, prompt-engineering evaluation, and clinical
guideline methodology.
\end{abstract}

\section{The problem: every prompt starts from zero}

Most AI writing assistance starts cold every time. There is no memory
of what worked last time, and no sense of an organisation's voice, its
standards, or its audience. The result is generic, confident output
that must be corrected after the fact.

The fix is not a better prompt. It is giving the model a documented
picture of what ``good'' looks like for a specific organisation, and a
way to improve that picture as more is learned about what works. Done
well, the model becomes more useful over time rather than requiring
the same corrective effort on every task.

A failure mode persists even once this is addressed: \emph{confident
wrongness with silent drift}. Output sounds authoritative but cites a
source incorrectly, or lands near the target word count while quietly
shifting the intended meaning. These errors are harder to catch than
obvious nonsense, and they compound faster.

\section{Supervised co-writing as the starting point}

Most people using AI for writing are not doing autonomous AI
publishing. They are doing \emph{supervised co-writing}: the model
handles a first pass or a tedious sub-task, a human reviews and
corrects, and a human publishes. That is the normal case, and it is
where most of the current value sits.

CADI targets teams who want to move from ``AI is something we
sometimes use'' toward ``AI is something our editorial standards
reliably shape.'' That transition does not happen through prompting
technique alone; it requires a documented, testable, and revisable
guidance layer.

\section{Preconditions}

The framework works best when an organisation can answer yes to the
following:

\begin{itemize}[nosep]
  \item There are ideally 50+ published items considered high quality
        (not merely recent). Ten to twenty items from a narrow,
        high-quality domain still work, at the cost of some
        overfitting; the discipline of running the loop is what
        compounds over time.
  \item Some quality signal exists: engagement data, editor picks,
        longevity tags, or manual curation.
  \item At least one person has time to own the guidance document and
        keep it current.
  \item Content can be exported from the CMS in text form (Markdown,
        HTML, or CSV).
  \item Editors are willing to log failures rather than silently
        fixing them.
\end{itemize}

None of this requires a data-science team, a custom CMS, or a formal
AI strategy. If no one will own the guidance document, however, the
method breaks down regardless of tooling.

\section{The CADI framework}

CADI is a four-phase loop, built once and then run continuously on its
final phase. Each cycle sharpens the guidance document rather than
adding to its weight.

\subsection{C --- Collect: extract the best examples}

An organisation's best content already embodies its patterns: word
counts, title structures, tag usage, narrative flow. The task is to
curate an exemplary set, not simply pull the most recent items.

Work within a single content subtype at a time. Most organisations
publish several superficially similar formats --- press releases,
news briefs, feature reports, donor updates, blog posts --- each with
its own voice and structure; mixing them yields average patterns that
fit none of them.

Top-traffic items are the obvious and most common wrong default.
High traffic often correlates with topic timeliness or reach rather
than the qualities worth reproducing. Three signals combine better:

\begin{itemize}[nosep]
  \item \textbf{Editorial investment} --- items the team spent the
        most time on, since effort is a quality signal even absent
        strong metrics.
  \item \textbf{Aspirational fit} --- items the organisation wants
        more of its output to resemble.
  \item \textbf{Performance for type} --- a press release performing
        at five times the median for press releases is more
        informative than a single feature that went viral.
\end{itemize}

Existing reference material --- style guides, brand guidelines,
glossaries, terminology lists, institutional memos --- should be
added alongside the corpus. The corpus shows the model what good looks
like by example; the reference material states the rules to enforce
and should be citable directly in generated drafts.

\emph{Exit criterion:} 50--100 items in one consolidated, cleaned
document; a written record of why they qualify; an editor's spot-check
for outliers.

\subsection{A --- Analyze: find quantitative patterns}

The curated examples are fed to a model with a structured prompt to
extract empirical description, not yet guidance:

\begin{itemize}[nosep]
  \item Average word counts, with ranges, by content type
  \item Title formats and syntactic shapes
  \item Heading and structural patterns
  \item Voice markers (person, contractions, hedging, sentence-length
        variance)
  \item Source treatment (inline links, blockquotes, first-paragraph
        attribution)
  \item Opening and closing patterns
  \item Cross-linking density
  \item Metadata usage
\end{itemize}

This list is not exhaustive; each content subtype has its own tells,
and the signals watched for should be documented so they can be
revisited on the next cycle. Models handle the extraction well; the
interpretation remains a human judgment.

\emph{Exit criterion:} a written summary of patterns with numbers,
reviewed by at least one editor for anything unrepresentative.

\subsection{D --- Document: turn patterns into working instructions}

The Analyze output is combined with existing institutional knowledge
into a single document the AI tool uses as standing context. Where an
existing style guide already covers a subtype, the Document step
tightens it with the Analyze step's numbers rather than re-deriving it
from scratch, overriding only where there is evidence the prior
guidance was wrong.

Observed patterns and existing rules will sometimes conflict --- for
example, an 8-word average title length against a house-style cap of
6. The decision, and its rationale, belongs in a change log rather
than remaining implicit.

The result is guidance concrete enough to act on (``180 words
$\pm$10\%'') rather than vague (``be concise''); the same document also
serves human editors.

\emph{Exit criterion:} the instructions reviewed by an editor not
involved in writing them, and integrated with the AI tool in use.

\subsection{I --- Iterate: build confidence over time}

Iteration done badly is tweaking in the dark. Done well, it is
hypothesis-driven, for example:

\begin{itemize}[nosep]
  \item \emph{If} explicit numeric targets are added (``180 words
        $\pm$10\%''), \emph{then} first-pass acceptance rises from 40\%
        to 70\% (illustrative).
  \item \emph{If} exemplars of title styles are shown, \emph{then} the
        model chooses appropriate titles 85\% of the time
        (illustrative).
  \item \emph{If} oppositional review is run (a second model critiques
        the guidance), \emph{then} biased rules are caught 80\% of the
        time (illustrative).
\end{itemize}

Hypotheses are tested with a simple before/after comparison on real
content, a concrete outcome metric, and a written record of what
changed and why. For a small team (two to three editors), 20--30 items
per condition supports an honest qualitative comparison --- an
improvement over blind tweaking, though not a statistical signal.
Genuine statistical rigor requires much larger, blinded samples that
most editorial teams will not reach; qualitative honesty about what
changed and why is the realistic bar.

A mid-cycle revision trigger applies when more than 20\% of a week's
outputs fail the same check: that pattern warrants investigation
immediately rather than waiting for a scheduled audit.

\emph{Exit criterion:} at least one completed hypothesis test, a
documented decision (kept or rejected), and a failure log someone is
actively adding to.

\section{Case study: CADI at UNDRR}

One full loop illustrates the framework in practice:

\begin{itemize}[nosep]
  \item \textbf{Collection}: 122 publications flagged as
        ``evergreen + high-engagement'' in Drupal.
  \item \textbf{Analysis}: the model extracted patterns --- 169-word
        average, 56\% of titles using colons, all items using two to
        four themes.
  \item \textbf{Documentation}: a guidance document was drafted and
        uploaded as standing instructions in Microsoft 365 Copilot.
  \item \textbf{Iteration}: quarterly audits checked whether new AI
        drafts matched the documented patterns, with a shared change
        log recording adjustments and rationale.
\end{itemize}

Before CADI, none of curation, pattern extraction, guidance
documentation, or iteration reliably produced anything usable in
isolation; each task restarted from nothing. Running them as a
connected loop converted a set of one-off attempts into a continuous,
improving process. The value was not a single headline metric but the
conversion of ad hoc effort into something repeatable. This case is
not offered as rigorous evidence: a prospective, quantified A/B test
of which guidance changes move which outcomes by how much remains the
missing piece, and is the direction future cycles should move toward.

\section{Where this breaks down}

Iteration reverts to guesswork under several conditions:

\begin{enumerate}[nosep]
  \item \textbf{Encoding bias without inspection} --- a pattern such
        as ``short sentences are good'' may simply reflect that the
        exemplars targeted beginners; applied universally, it
        infantilizes advanced content.
  \item \textbf{Optimizing for fit rather than truth} --- ``this tag
        co-occurs 26\% of the time, so always pair them'' should read
        ``consider pairing them.''
  \item \textbf{No quality gate on the source set} --- mediocre
        content in the ``best'' set skews guidance toward mediocrity.
  \item \textbf{Stopping too early or too late} --- ten examples is
        too few; five hundred risks overfitting.
  \item \textbf{Expecting voice cloning in informal genres} ---
        style-imitation performs well for structured formats such as
        news and email, and considerably worse for informal, personal
        writing such as blogs and forum posts; corpus-derived guidance
        reliably reproduces structure and convention, but a
        distinctive personal voice still requires a human pass.
\end{enumerate}

Mitigations include transparent curation criteria, documented
assumptions, human review between cycles, oppositional analysis (a
second model critiquing the instructions), and explicit escape hatches
(``usually X, unless Y'').

CADI also does not fit cleanly for high-volume newsrooms shipping
dozens of pieces per week across many content subtypes; it assumes one
primary subtype at a time and a corpus small enough to read in full.
Sampling methodology and per-subtype parallel runs would be needed at
that scale, and are out of scope here.

\section{Related work}

None of the individual pieces of CADI are novel in isolation; the
contribution is the connected loop.

Commercial tools such as Writer and Jasper reverse-engineer a voice
profile from an organisation's own content, functioning as
implementations of the Collect and Analyze phases. ChainForge (CHI
2024) formalized treating prompts, and the criteria used to evaluate
them, as testable hypotheses --- the same posture CADI takes toward
its own guidance document. Clinical guideline methodology (AGREE II
and GRADE) assesses whether guidelines are evidence-based and
implementable; the Australian COVID-19 living guidelines went further,
revising and republishing weekly for 24 weeks as new evidence arrived
--- a continuously iterated guidance document running in production.

A software subscription is not a method, and a research paper is not a
guide. What appears missing elsewhere is the connected discipline:
deriving guidance from a curated, organisation-specific corpus, then
maintaining the failure log and test record that make that guidance
revisable on evidence rather than intuition. That combination is what
CADI names.

\section{Conclusion}

CADI is best understood as a forcing function rather than a generator:
most of the underlying work is still careful reading of an
organisation's own content. The framework's contribution is ensuring
that reading happens, and that what is learned from it gets written
down, tested, and revised. It is presented here as the method run
in production at UNDRR; adaptation should be expected for different
volumes, content types, and organisational constraints.

\section*{References}

\begin{enumerate}[nosep]
  \item Bsharat et al., ``Principled Instructions Are All You Need,''
        \href{https://arxiv.org/abs/2312.16171}{arXiv:2312.16171}.
  \item ``Catch Me If You Can? Not Yet: LLMs Still Struggle to Imitate
        the Implicit Writing Styles of Everyday Authors,'' EMNLP 2025
        Findings, \href{https://arxiv.org/abs/2509.14543}{arXiv:2509.14543}.
  \item Wu et al., ``ScatterShot: Interactive In-context Example
        Curation for Text Transformation,'' ACM IUI 2023,
        \href{https://arxiv.org/abs/2302.07346}{arXiv:2302.07346}.
  \item Arawjo et al., ``ChainForge: A Visual Toolkit for Prompt
        Engineering and LLM Hypothesis Testing,'' CHI 2024,
        \href{https://arxiv.org/abs/2309.09128}{arXiv:2309.09128}.
  \item Tendal et al., ``Weekly updates of national living
        evidence-based guidelines: methods for the Australian living
        guidelines for care of people with COVID-19,'' \emph{Journal
        of Clinical Epidemiology}, 2020,
        \href{https://pmc.ncbi.nlm.nih.gov/articles/PMC7657075/}{PMC7657075}.
  \item Wei et al., ``Chain-of-Thought Prompting Elicits Reasoning,''
        2022, \href{https://arxiv.org/abs/2201.11903}{arXiv:2201.11903}.
  \item Yao et al., ``Tree of Thoughts: Deliberate Problem Solving with
        LLMs,'' 2023, \href{https://arxiv.org/abs/2305.10601}{arXiv:2305.10601}.
  \item Kohavi, Longbotham et al., ``Controlled experiments on the
        web: survey and practical guide,'' \emph{Data Mining and
        Knowledge Discovery}, 2009,
        \href{https://doi.org/10.1007/s10618-008-0114-1}{doi:10.1007/s10618-008-0114-1}.
  \item AGREE II --- \href{https://www.agreetrust.org/}{agreetrust.org}.
  \item GRADE --- \href{https://www.gradeworkinggroup.org/}{gradeworkinggroup.org}.
  \item Di\'ataxis --- \href{https://diataxis.fr/}{diataxis.fr}.
  \item PRISMA --- \href{https://www.prisma-statement.org/}{prisma-statement.org}.
\end{enumerate}

\end{document}
